\medskip
\hypertarget{native-108-article079}{}\pdfbookmark[2]{79.}{article-079}
\begin{center}79.\end{center}
Zur vollständigen Bestimmung eines Kegelschnitts in seiner Ebene braucht man \emph{drei} Größen, nämlich die Lage des Perihels, die Exzentrizität und den Halbparameter. Sind diese aus gegebenen, von ihnen abhängigen Größen zu ermitteln, so müssen genügend Daten vorliegen, um drei voneinander unabhängige Gleichungen aufstellen zu können. Jeder nach Betrag und Lage gegebene Radiusvektor liefert eine Gleichung. Zur Bestimmung der Bahn braucht man daher drei nach Betrag und Lage gegebene Radiusvektoren. Liegen nur zwei vor, so muss entweder eines der Elemente selbst schon gegeben sein oder wenigstens eine andere Größe, aus der sich eine dritte Gleichung aufstellen lässt. Daraus ergeben sich verschiedene Aufgaben, die wir nun der Reihe nach behandeln.
\par
Es seien $r$, $r'$ zwei Radiusvektoren, die mit einer beliebigen, von der Sonne aus in der Bahnebene gezogenen Geraden die Winkel $N$, $N'$ bilden, in Bewegungsrichtung gezählt. Ferner sei \SourceLeaf{109}$\Pi$ der Winkel, den der Radiusvektor im Perihel mit derselben Geraden bildet, sodass den Radiusvektoren $r$, $r'$ die wahren Anomalien $N-\Pi$, $N'-\Pi$ entsprechen. Schließlich sei $e$ die Exzentrizität und $p$ der Halbparameter. Dann gelten die Gleichungen
\[
\begin{aligned}
\frac{p}{r}&=1+e\cos(N-\Pi)\\
\frac{p}{r'}&=1+e\cos(N'-\Pi),
\end{aligned}
\]
aus denen sich die beiden übrigen Größen bestimmen lassen, wenn außerdem eine der Größen $p$, $e$, $\Pi$ gegeben ist.
\par
Nehmen wir zunächst an, der Halbparameter $p$ sei gegeben. Dann lassen sich die Größen $e$ und $\Pi$ offenbar aus den Gleichungen
\[
\begin{aligned}
e\cos(N-\Pi)&=\frac{p}{r}-1\\
e\cos(N'-\Pi)&=\frac{p}{r'}-1
\end{aligned}
\]
nach Lemma III des vorangehenden Artikels bestimmen. Somit gilt
\[
\begin{aligned}
\operatorname{tang}(N-\Pi)&=\operatorname{cotang}(N'-N)-\frac{r(p-r')}{r'(p-r)\sin(N'-N)}\\[5pt]
\operatorname{tang}(\tfrac12N+\tfrac12N'-\Pi)&=\frac{(r'-r)\operatorname{cotang}\tfrac12(N'-N)}{r'+r-\dfrac{2rr'}{p}}\mathpunct{\cdot}
\end{aligned}
\]
